\chapter{Discussion and Conclusion}
\label{chapter:discussion}

\section{Review of the results}
\label{sec:results-review}

The construction of Chapter~\ref{chapter:methodology} rests on a structural claim: REINFORCE is biased in mean-field control because it differentiates the objective at a frozen population flow, and the missing term can be recovered model-free by randomizing the population argument. The experiments support that claim, and also show that recovering the missing term is not by itself a reason to pay for it.

The suite does not order the estimators uniformly, and that disagreement is the most informative feature of Table~\ref{tab:objective-summary}. Three regimes appear. On the two-state problem at $T=5$, on the linear--quadratic benchmark and on the portfolio benchmark, an estimator that differentiates the population flow wins by a margin far larger than the seed dispersion, an order of magnitude on two-state. On distribution planning and targeted advertising the population-flow derivative helps, but the logit-perturbed estimator extracts more of it than transport does. On cybersecurity plain REINFORCE is the best method, and the entire spread between best and worst is under $6\%$ of the objective, so that benchmark separates the methods weakly in absolute terms. One entry deserves separate comment: on distribution planning the adaptive controller is not merely worse than the fixed grid but worse than every other method.

Matched simulator budgets, moreover, do not imply matched cost. The transport update spends its budget on an auxiliary sensitivity estimate that is cheap in simulated transitions and expensive in arithmetic, so its wall-clock cost exceeds that of REINFORCE by a factor between $2$ and $26$. The ratio is a property of the dimension of the population argument rather than of the estimator: the extreme case is distribution planning, whose nine-coordinate sensitivity makes the recursion the dominant cost of the update and which is also the benchmark with the smallest auxiliary sample size relative to that dimension.

The clearest positive evidence is the two-state benchmark at $T=5$. REINFORCE degrades over iterations: it ends at $-31.51$ against an initialization of $-4.53$ and an optimum of $-2.64$, while the transport estimator reaches $-3.13$ and recovers the closed-form optimal policy to within $3.5\%$. The same pattern, with smaller margins, appears on both continuous-state benchmarks, where transport closes three quarters of the optimality gap on the linear--quadratic problem and beats REINFORCE by more than twice the seed dispersion on the portfolio problem. In all three cases the population channel carries the sign of the useful direction, and an estimator that ignores it optimizes a different function.

The counterexamples are equally informative: on cybersecurity plain REINFORCE is the best method at a fraction of the wall-clock cost, and on distribution planning and targeted advertising the logit-perturbed estimator of \cite{meunier2026model} extracts more of the population-flow derivative than transport does. Section~\ref{sec:correction-alignment} isolates what separates the regimes, and it is not the magnitude of the correction but its alignment $\varrho$ with the policy-score term. The portfolio ablation makes this controlled rather than correlational: switching the market-impact coefficient $\gamma$ from $2$ to $0$ leaves the dynamics unchanged, drives $\varrho$ from $-0.196$ to $+1.000$, and reverses the ordering of the estimators.

Section~\ref{sec:theory-validation} adds the check that the estimators behave as the analysis says. All four testable predictions hold in the regime the theory targets, and the bias--variance trade-off in $\lambda$ is therefore not a formal artefact of the bounds but a measurable feature of the estimator, with an interior optimum on every benchmark where a grid was run.

\section{Strengths and limitations}
\label{sec:strengths-limitations}

\paragraph{Strengths.}
Four properties distinguish the estimator from the finite-state logit perturbation it generalizes. It is stated at the level of probability measures, so the same argument produces a simplex perturbation in finite state spaces and a transport perturbation in continuous ones, the representation entering only through the map $\Psi$. The simplex version does not require strictly positive population coordinates, which the logit chart does, so it remains defined on the boundary of the simplex where several of the learned flows actually go. The forward recursion for the sensitivity makes the simulator cost linear in the horizon rather than quadratic, and the two-state experiment isolates exactly that effect: at $T=2$ the logit estimator is the more accurate, and at $T=5$ the ordering reverses, its policy error growing by a factor $35$ against $2.5$ for transport at the same budget. Finally, because the correction is an explicit term, its usefulness can be tested before committing to it; the alignment diagnostic costs two automatic-differentiation passes and predicts, without a training grid, whether REINFORCE will be directionally competitive.

\paragraph{Limitations of the method.}
The estimator is not uniformly better, and the reasons are worth stating plainly rather than as a caveat. Matching the simulator budget does not match the wall-clock cost, since the transport update carries the auxiliary recursion and its automatic-differentiation graph, so where REINFORCE is directionally correct it is not merely competitive but strictly preferable. The correction is also worth its variance only when the auxiliary sensitivity can be resolved: distribution planning allocates $1.2$ auxiliary trajectories per free coordinate of a nine-dimensional population argument, the worst ratio in the suite, and is where transport loses most heavily to the logit estimator; the portfolio benchmark is the same phenomenon at the other extreme, transport winning only once the auxiliary sample size is raised from $1$ to $200$. The adaptive controller is the weakest component: it contracts the main scale below the smallest calibrated value on every finite-state benchmark and is the worst method on distribution planning, because at four replications per checkpoint the debiased bias estimate is frequently unresolved, and reading that as an absence of bias drives the controller towards a variance-minimizing corner.

\paragraph{Limitations of the evidence.}
Other conclusions are bounded by the design of the study rather than by the method. The continuous-state theory is developed only as far as objective-level convergence, so gradient-level convergence, the policy-gradient formula, the sensitivity estimator and the error decomposition remain open in the Gaussian-mixture case, and the continuous-state experiments are run with an estimator whose guarantees are an extrapolation. The replicated gradient diagnostics need an exact oracle and a directly parametrized policy, which only two benchmarks provide, and on the portfolio benchmark the unbiasedness test is inconclusive rather than passed, since at a signal-to-noise ratio of $0.02$ per replication fifty replications cannot resolve a bias of the size the theory allows. The auxiliary error bound of Proposition~\ref{prop:auxiliary-sensitivity-error} is the one prediction the saved artifacts cannot separate into its two terms. Finally, five seeds resolve the comparisons between estimator families by wide margins but not the ordering within the transport grid on the portfolio benchmark, where the seed dispersion exceeds the spread of the grid.

\section{Extensions}
\label{sec:extensions}

\paragraph{An actor--critic form of the estimator.}
The most direct improvement the analysis suggests is variance reduction on the return. In \eqref{eq:discrete-transport-estimator} the perturbation score is multiplied by the full perturbed return and carries the factor $\lambda^{-1}$ that the mean-square-error decomposition identifies as the dominant variance term at small scales; a deterministic baseline removes only a constant, whereas an advantage estimated from a learned critic would remove the part of the return uninformative about the current step. The critic is a natural object here, since the perturbed problem is an ordinary Markov decision process on the augmented space in which the randomized population coordinate is part of the state, as the augmented trajectory law \eqref{eq:gmm-augmented-trajectory-law} of Appendix~\ref{appendix:gmm-augmented} makes explicit, so $V^\theta(t,x,m)$ is well defined. Two questions must be answered first. The critic is evaluated at the perturbed population argument, so its approximation error enters through the same $\lambda^{-1}$ prefactor as the score, and the bias it introduces has to be traded against the variance it removes in Theorem~\ref{theorem:discrete-mse-decomposition}. The sensitivity recursion also has to be revisited, since the critic depends on $\theta$ through the population flow exactly as the reward does.

\paragraph{Beyond exchangeability.}
The construction rests on exchangeability: a single representative agent, one population law, one perturbation per step. Two relaxations are natural. In non-exchangeable models the agents are coupled through a graph rather than a single mean field and the limit is described by a graphon, so the population argument becomes a function-valued object indexed by a continuum of vertex types. The framework does not obviously break there, since it requires only a finite-dimensional representation of the population argument and a probability-preserving randomization of it, and a graphon admits finite-dimensional approximations of exactly the kind $\Psi$ is designed to accommodate; what must be redone is the sensitivity recursion, which currently propagates a single flow rather than a family indexed by type. Finitely many interacting populations, where the argument is a tuple of laws, is the natural first step and needs little additional theory.

The second relaxation is extended mean-field control, where the coefficients depend on the joint law of the state and the action. This is the more immediate of the two, because the estimator is indifferent to what the population argument represents: it needs a finite-dimensional chart, a randomization within it, and a forward recursion for its sensitivity, and the joint law of $(X_t,\alpha_t)$ has all three whenever the action space is finite or the joint law is summarized by moments. The difficulty is that the policy then influences the population argument twice, directly through the sampled actions and indirectly through the state flow, so the decomposition into a policy score and a population score acquires a third contribution the present derivation does not have.

\paragraph{Completing the continuous-state analysis.}
The gap identified above is the natural theoretical continuation: gradient-level convergence for the Gaussian-mixture perturbation, the corresponding policy-gradient representation, a model-free estimator of the sensitivity $D_\theta z_{K,t}^\theta$ of the representation coordinates, and the error decomposition. The finite-state proofs give the template, and the additional difficulty is the representation error $\delta_K^\theta$, which does not vanish with $\lambda$ and must be carried through the differentiated estimates. A related and more open question is whether $\Psi$ can be learned rather than fixed, since the construction requires only that the representation be finite-dimensional and that its coordinates be differentiable in $\theta$.

\paragraph{Allocating the auxiliary budget.}
The equal-budget rule fixes the total number of simulated transitions per update but not their split between the main trajectories and the auxiliary sensitivity estimate, and that split, rather than the estimator itself, is what limits transport on the higher-dimensional simplex. A systematic study of the allocation as a function of the dimension of the population argument and of the horizon is the most direct extension of the present results. The diagnostics of Section~\ref{sec:theory-validation} give that the auxiliary sample size enters through the $(n\eta^2)^{-1}$ term of Proposition~\ref{prop:auxiliary-sensitivity-error}, which is measurable at a frozen policy without running a training grid.

\paragraph{Constraining the adaptive controller.}
Two modifications follow directly from the diagnosis above: a lower bound on $\lambda$ at the smallest scale for which the perturbation-density score has been validated, and a balance target that accounts for the dimension of the simplex, since the score has $N-1$ coordinates and scales as $\lambda^{-1}$. A third, less local change would be to resolve the bias estimate properly by accumulating it across checkpoints rather than re-estimating it from four replications each time.

\paragraph{A benchmark with a finite-dimensional law statistic.}
Both continuous-state benchmarks are driven by the population mean, so the representation is one-dimensional. The controlled noisy Kuramoto model \cite{sinigaglia_braghin_berman_kuramoto,chen_wang_zhidkova_kuramoto}, in which oscillators on the circle are steered towards synchronization, is a natural next case: although the state is an angle and the law a distribution on the circle, the interaction enters only through the first Fourier mode $(\E[\cos X_t],\E[\sin X_t])$, so the population argument is again finite-dimensional, of dimension two rather than one, and it would test whether the auxiliary recursion degrades when the law statistic gains a coordinate.

\section{Conclusion}

This report set out to determine whether policy gradients for discrete-time mean-field control can be made fully model-free without the finite-state logit representation on which the existing construction depends. The answer developed here is a perturbation framework stated at the level of probability measures: the population argument is replaced by the pushforward of a random transport map, reducing to a probability-preserving mixture on the simplex in finite state spaces and to a random displacement read through a finite-dimensional representation in continuous ones. In both cases the gradient of the perturbed objective splits into the familiar policy score and a population score, both estimable from simulated trajectories without differentiating the transition kernel, the reward, or the population recursion, and a forward recursion for the sensitivity flow makes the simulator cost linear rather than quadratic in the horizon. The two settings are not at the same stage, and Table~\ref{tab:contribution-status} states the difference plainly: on a finite state space the argument is complete, from the perturbation estimate through gradient-level convergence to the mean-square-error decomposition, whereas on a continuous state space it is established at the level of the objective and open at the level of the gradient. The continuous-state estimator is therefore built by analogy with the finite-state one and supported by measurement rather than by proof.

Empirically the picture is more interesting than a uniform improvement would have been. Across six benchmarks at matched simulator budgets no estimator dominates, and the disagreement is itself the main result: whether REINFORCE is directionally competitive is governed by how far the population-flow correction is from collinear with the term it already computes, a condition that a cheap evaluation-only diagnostic detects in advance and that a controlled ablation confirms. Where the correction is strongly non-collinear, population-aware estimators beat REINFORCE outright, by an order of magnitude on the two-state problem and by a wide margin on the portfolio benchmark. Where it vanishes, as on cybersecurity, REINFORCE is the better choice and the perturbation only adds variance. The diagnostic is silent about what happens between those poles, and the linear--quadratic benchmark is exactly that case: the alignment is above $0.98$, so REINFORCE is directionally right and captures most of the available improvement, yet a small orthogonal component leaves it a visible residual gap, three quarters of which transport closes.

What remains is set out above: completing the gradient-level analysis in the continuous-state case, reducing the estimator variance through a critic, allocating the auxiliary budget in a principled way, and relaxing exchangeability. The internship continues until the end of October, and these are the directions in which the work is being taken.
